\documentclass[12pt,a4paper]{article}

% ====== PACKAGES ======
\usepackage[utf8]{inputenc}
\usepackage{geometry}
\usepackage{graphicx}
\usepackage{titlesec}
\usepackage{enumitem}
\usepackage{setspace}
\usepackage{hyperref}
\usepackage{authblk}
\usepackage{booktabs}
\usepackage[backend=biber,style=numeric,citestyle=numeric]{biblatex}

\addbibresource{sources.bib}
% ====== PAGE SETTINGS ======
\geometry{margin=1in}
\setstretch{1.3}

\makeatletter
\renewcommand{\maketitle}{
	\begin{center}
		{\LARGE \bfseries \@title \par}
		\vskip 1em
		{\large \@author \par}
		\vskip 0.5em
		{\@date}
	\end{center}
}
\makeatother

% ====== SECTION STYLE ======
\titleformat{\section}{\large\bfseries}{\thesection.}{1em}{}
\titleformat{\subsection}{\normalsize\bfseries}{\thesubsection}{1em}{}

% ====== TITLE INFO ======
\title{\textbf{Ontologically Faithful Generation of Non-Player Character Dialogues
	}\\[0.5em]
}
\author[1]{Sina Taheri Behrooz}
\author[2]{Hesam soleymani}
\affil[1]{Iran University of Science and Technology (IUST), sina\_taheri@comp.iust.ac.ir}
\affil[2]{Iran University of Science and Technology (IUST), email2@example.com}

\date{\today}

\begin{document}
\maketitle
\thispagestyle{empty}
\newpage

\setcounter{page}{1}

\section{Introduction}

One of the most important factors in creating believable and memorable characters in RPG games is the character dialogue. Dialogue serves as the primary medium through which players perceive character personalities, motivations, and emotions, shaping the overall narrative experience \cite{loyall1997believable}. However, Writing satisfactory dialogue requires time and effort from the game writers which in turn increases the development cost of games. Other than the cost and effort, the tree like structure of dialogue means because of the dependency between different nodes, a small change in the story creates a cascade of changes in the dialogue tree.

Early attempts to address these challenges utilized procedural methods, such as finite-state machines (FSMs) and tree-based discourse generation frameworks \cite{strong2008npc}. While these methods provided structured control over dialogue flow, they generally failed to produce contextually rich or emotionally coherent conversations.

With the emergence of large language models (LLMs), new research has been done on the  creation of more natural, consistent, and context-aware character dialogue. Recent papers such as KNUDGE \cite{weir2024knudge} demonstrates that ontology-grounded dialogue generation can improve narrative consistency and factual correctness, while systems like LIGHT \cite{urbanek2019light} show that grounding dialogue in a simulated game world enhances coherence and immersion.

Building upon these developments, this project aims to explore ontology-informed generative approaches for automating NPC dialogue creation, reducing authoring costs while preserving believability and narrative consistency.

\section{Problem Statement}
The KNUDGE paper utilizes LLMs as the main method for graph generation, yet LLMs are inherently optimized for producing linear text. Consequently LLMs are prone to generating incoherent and inconsistent dialogue trees which in turn increase the time required by writers and designers to flesh out unique and immersive dialogue. 

This project investigates whether applying structured graph serialization techniques, prompt engineering strategies and utilizing modern LLMs with enhanced reasoning and expanded context windows can improve the coherence, structural validity, and ontology fidelity of the dialogue trees generated by the base KNUDGE system.

\section{Data Sources}
The main data source for developing this project is the KNUDGE \cite{weir2024knudge} Dataset. This dataset contains 159 dialogue trees from 45 side quests of \textit{The Outer Worlds}. In this dataset, each dialogue has been tagged and annotated to form a tree like structure. Also specific interactions where progression is tied to character attributes have been marked via watching playthroughs. This dataset also contains support knowledge related to each entity required for generating dialogues. This support knowledge has been scraped from fan wikis related to the game.

\section{Aims}
Our aim is to conduct an ablation study on effects of the following methods on the core evaluation metrics of KNUDGE: 
	\begin{itemize}
		\item To evaluate the effect of structured graph serialization techniques on LLMs creating coherent and well formed dialogue trees.  
		\item To assess the effectiveness of prompt engineering methods in enhancing ontology usage. 
		\item To examine whether modern LLMs with expanded context window and reasoning outperform the models used in base paper.
	\end{itemize}

\section{Methodology}

The goal of this project is to evaluate how graph serialization, prompt engineering, and modern LLMs influence the quality of dialogue-tree generation within the KNUDGE framework. All experiments will use the same base dataset and evaluation pipeline as the original KNUDGE paper to ensure comparability. 
\subsection{Graph Serialization}
the base paper uses a JSON based graph serialization method. In this format, critical information about the graph such as the start and end nodes, general data about the quest and an edge list representing the dependencies is stored. We propose 3 representation methods for this data. After exemplars are retrieved using the BM-25 index, their representation will be changed based on the method we are using before being injected into the prompt. 
\begin{itemize}
	\item \textbf{Adjacency List} : Each node explicitly lists all of it's neighboring nodes. 
	\item \textbf{Adjacency Matrix} : A 2d matrix is replaced with the edge list.
	\item \textbf{GML} : Represent the whole document in a new graph specific format. Attributes of each dialogue will be inside GML node while each edge shows the relation between 2 nodes.  
\end{itemize}
These formats allow us to test whether more structured or more explicit graph encodings improve the structural fidelity and ontology adherence of the generated graphs.

\subsection{Prompt Engineering}
There has been some research lately on improving an LLM's graph generation ability utilizing prompt engineering techniques \cite{demirci-etal-2025-llms}\cite{yao2024exploringpotentiallargelanguage}. The KNUDGE paper utilizes iterative prompting as a technique for generating end-to-end dialogue trees, we plan to study the effects of 2 other techniques :
\begin{itemize}
	\item \textbf{Direct prompting} : The model receives a single comprehensive prompt containing the entire task description and few examples. No further iteration is done on the result. 
	\item \textbf{Program-Augmented} Prompting : The model receives the verification script as another part of the prompt and is asked to refer to the verification script at the time of generation.
\end{itemize}
\subsection{Modern LLMs}
The KNUDGE paper uses the gpt-4-turbo\cite{openai2024gpt4technicalreport} and Vicuna-13B-v1.5-16k models for the end-to-end dialogue write and the older text-davinci-003 GPT-3 model for the node dialogue writer. We will utilize 3 newer models with enhanced reasoning ability and larger context. 
\begin{itemize}
    \item \textbf{Gemini 2.5 Pro}: A modern reasoning-optimized model with large context capacity.
	\item \textbf{Qwen2 Family}: A suite of open models with various sizes, enabling experimentation under different hardware constraints.
	\item \textbf{Llama 4 Family}: The most recent iteration of the Llama series.
\end{itemize}
The specific parameter sizes of Qwen2 and Llama models will be chosen based on available compute resources, prioritizing variants that can run locally without extensive GPU requirements.


\section{Conclusion}
This project proposes a systematic investigation into how graph serialization, prompt engineering, and modern large language models influence the generation of dialogue trees within the KNUDGE framework. Together, these experiments aim to provide a clearer understanding of the factors that most significantly impact controlled dialogue-tree generation. The findings will offer practical insights for improving narrative generation pipelines, guiding future LLM selection, and informing best practices for structured content generation in interactive storytelling environments.

\printbibliography[title={References}]

\end{document}
